% 01-introduction.tex -- Introduction and Program Overview
\chapter{Introduction and Program Overview}
\label{ch:introduction}

% ====================================================================
\section{About the Program}
\label{sec:about}

The MSCA Industrial Doctoral Network on Digital Finance---\emph{Reaching New
Frontiers}---is a prestigious research and training program funded under
Horizon Europe (Grant Agreement No.~101119635) with a total EU contribution of
approximately EUR~4~million.  Running from January~2024 through December~2027, the network
brings together nineteen partners from eleven European countries to train a
new generation of researchers at the intersection of finance, artificial
intelligence, and data science.

Seventeen doctoral candidates (DCs) conduct their research across five thematic
work packages, supported by a comprehensive curriculum spanning mandatory
courses, advanced electives, and transferable-skill workshops.  Each DC benefits from joint supervision by academic and industry
mentors, international secondments, and exposure to regulatory and policy
institutions.

As the program enters its third year, 2026 marks a period of intensified
research output, expanded training activities, and---for the first time---broad
external participation.  The consortium has designed the 2026 training calendar
to be accessible to PhD students across Europe, regardless of whether they hold
one of the seventeen funded positions.

\vspace{8pt}
\begin{center}
  \numbercallout{17}{PhD Positions}
  \hfill
  \numbercallout{11}{Countries}
  \hfill
  \numbercallout{19}{Partners}
  \hfill
  \numbercallout{40}{Courses}
\end{center}
\vspace{8pt}

\placeholderphoto{0.8\textwidth}{European research campus}

% ====================================================================
\section{Open Training Philosophy}
\label{sec:open-training}

The Digital Finance network is committed to an open training philosophy.  While
the seventeen funded positions are the core of the program, the majority of
training activities---summer schools, workshops, lecture series, and
seminars---are open to \textbf{all European PhD students} working in digital
finance and related fields.  Most events are offered free of charge; where
capacity is limited, registration is allocated on a first-come, first-served
basis.

Researchers from outside the consortium are warmly encouraged to apply.
Detailed registration instructions, deadlines, and event descriptions are
available on the program website.

\begin{highlightbox}
\textbf{How to Join}\\[4pt]
Visit the program website for registration details, upcoming events, and
application deadlines.  Scan the QR code below or go to
\url{https://www.digital-finance-msca.com/}.

\vspace{6pt}
\begin{center}
  \qrcode[height=2cm]{https://www.digital-finance-msca.com/}
\end{center}
\end{highlightbox}

% ====================================================================
\section{How to Use This Booklet}
\label{sec:how-to-use}

This booklet serves as the definitive guide to the 2026 training program.
Chapter~\ref{ch:consortium} introduces the consortium and its partners.
Chapter~\ref{ch:research-wps} outlines the five research areas.
Chapters~\ref{ch:training-overview} through~\ref{ch:transferable-skills} detail
the training structure, mandatory courses, elective offerings, and transferable
skills.  Chapter~\ref{ch:secondments} explains the secondment program.  The
remaining chapters cover scheduled events, participation guidelines, and
practical information.
